\chapter{کدگذاری موقعیت}

\section{جایگاه این بحث در درس}
پس از پایان بحث \eng{Embedding} و معرفی معماری \eng{Sequence-to-Sequence} برای \eng{G2P}، پیش از ادامهٔ آن معماری باید یکی از مشکلات مهم \eng{Embedding} پویا (بافتی) و راه‌حل آن را بشناسیم: \term{کدگذاری موقعیت}{Positional Encoding}. مفاهیم \eng{Query}، \eng{Key}، \eng{Value} و \eng{Self-Attention} را در فصل قبل دیدیم و اینجا فقط مرور می‌کنیم.

\begin{exambox}{نکتهٔ روی تخته}
کدگذاری موقعیت \emph{فقط در \eng{Embedding} پویا معنا پیدا می‌کند}؛ \eng{Embedding} ایستا مستقل از بافت است و مسئلهٔ «نقش کلمه نسبت به کلمه‌های دیگر» در آن مطرح نیست.
\end{exambox}

\section{مسئلهٔ ترتیب در \eng{Attention}}

دو جملهٔ «ایران آمریکا را شکست داد» و «آمریکا ایران را شکست داد» کلمه‌های یکسان ولی معنای کاملاً متفاوتی دارند. نمایش بافتی «شکست» در این دو جمله باید متفاوت باشد. در \eng{Self-Attention} کلمهٔ هدف نقش \eng{Query} را دارد، دیگر کلمه‌ها \eng{Key} و \eng{Value} دارند، وزن‌های توجه $\alpha_i$ از مقایسهٔ \eng{Query} با \eng{Key}ها ساخته می‌شوند و بردار بافت ترکیب وزن‌دار \eng{Value}هاست:
\[
\text{\lr{Context Vector}} \;=\; \sum_i \alpha_i V_i .
\]

\begin{figure}[htbp]
\centering
\resizebox{0.97\textwidth}{!}{%
\begin{tikzpicture}[rtl, 
  w/.style={font=\Large,inner sep=3pt},
  vec/.style={font=\large,text=mainblue},
  hl/.style={draw=red!70!black,thick,ellipse,inner sep=3pt}]
  % sentence 1 (right to left)
  \node[w] at (13.0,0) {\RL{جملهٔ ۱}};
  \node[w] (a1) at (10.2,0) {\RL{ایران}};
  \node[w] (a2) at (7.8,0) {\RL{آمریکا}};
  \node[w] at (5.6,0) {\RL{را}};
  \node[w,hl] (a4) at (3.4,0) {\RL{شکست}};
  \node[w] at (1.0,0) {\RL{داد}};
  \node[vec] at (10.2,-1.3) {$k_1=\binom{2}{0}$};
  \node[vec] at (7.8,-1.3) {$k_2=\binom{0}{2}$};
  \node[vec] at (3.4,1.4) {$Q=\binom{1}{1}$};
  \draw[-{Stealth},mainblue] (3.4,0.9)--(3.4,0.55);
  % sentence 2
  \node[w] at (13.0,-4.6) {\RL{جملهٔ ۲}};
  \node[w] (b1) at (10.2,-4.6) {\RL{آمریکا}};
  \node[w] (b2) at (7.8,-4.6) {\RL{ایران}};
  \node[w] at (5.6,-4.6) {\RL{را}};
  \node[w,hl] (b4) at (3.4,-4.6) {\RL{شکست}};
  \node[w] at (1.0,-4.6) {\RL{داد}};
  \node[vec] at (10.2,-5.9) {$k_2=\binom{0}{2}$};
  \node[vec] at (7.8,-5.9) {$k_1=\binom{2}{0}$};
  \node[vec] at (3.4,-3.2) {$Q=\binom{1}{1}$};
  \draw[-{Stealth},mainblue] (3.4,-3.7)--(3.4,-4.05);
  % scores
  \node[vec,align=center] at (3.4,-1.5) {$Q\!\cdot\! k_1=2,\ Q\!\cdot\! k_2=2$\\[2pt]$\alpha_1=\alpha_2=0{,}5$};
  \node[vec,align=center] at (3.4,-6.1) {$Q\!\cdot\! k_1=2,\ Q\!\cdot\! k_2=2$\\[2pt]$\alpha_1=\alpha_2=0{,}5$};
  \draw[gray!60,dashed] (-0.8,-2.7)--(14.6,-2.7);
\end{tikzpicture}}
\caption{مسئلهٔ ترتیب: کلمهٔ «شکست» (دایره‌شده) در هر دو جمله همان $Q=(1,1)$ را دارد و \eng{Key} کلمه‌های دیگر هم یکسان است؛ فقط جای آن‌ها عوض شده است. در هر دو جمله $Q\cdot k_1=Q\cdot k_2=2$ و وزن‌ها برابر می‌شوند، پس بدون اطلاعات موقعیت مدل تفاوتی نمی‌بیند. (بازسازی‌شده از تختهٔ کلاس)}
\end{figure}

\paragraph{مثال عددی کلاس.} با $q=(1,1)$ و $k_1=(2,0)$، $k_2=(0,2)$ داریم $q\cdot k_1=q\cdot k_2=2$، پس $\alpha_1=\alpha_2=0{,}5$ و اگر \eng{Value} هر دو ۱ باشد، $C=0{,}5\times1+0{,}5\times1=1$. با جابه‌جایی دو کلمه، \eng{Key} و \eng{Value} هر کلمه عوض نمی‌شود و جمع وزن‌دار دقیقاً همان ۱ را می‌دهد؛ یعنی مدل نمی‌تواند بفهمد کدام کلمه زودتر آمده است.

پرسش استاد: در دو جمله چه چیزی تغییر کرده است؟ پاسخ: \emph{موقعیت} کلمه‌ها. در رویکرد موازی مدل فقط «مجموعه‌ای» از نمایش کلمه‌ها را می‌بیند و ترتیب را صریحاً نمی‌داند. راه‌حل، افزودن اطلاعات موقعیت به نمایش هر کلمه است.

\begin{definitionbox}{کدگذاری موقعیت (\eng{Positional Encoding})}
افزودن اطلاعات «این عنصر چندمین عنصر دنباله است» به نمایش هر عنصر، تا مدل علاوه بر اینکه بداند هر کلمه \emph{چیست}، بداند \emph{کجا} قرار گرفته است. مثلاً در «ایران آمریکا را شکست داد»: «ایران» موقعیت ۱، «آمریکا» ۲، «را» ۳ و «شکست» ۴.
\end{definitionbox}

\begin{correctionbox}{دقت: «بی‌توجه به ترتیب» یعنی چه؟}
\eng{Self-Attention} بدون کدگذاری موقعیت \term{هم‌متغیر نسبت به جایگشت}{Permutation-Equivariant} است: بردار بافت یک \eng{Query} نسبت به جابه‌جایی فهرست \eng{Key}/\eng{Value}ها ثابت می‌ماند. این به معنای «\eng{Attention} بافت را نمی‌فهمد» نیست؛ کلمه‌ها محتوای هم را می‌بینند و تنها اطلاعات \emph{ترتیب} از دست می‌رود.
\end{correctionbox}

\section{آیا مدل‌های بازگشتی هم به کدگذاری موقعیت نیاز دارند؟}
\label{sec:ch10-rnn}
\begin{assignmentbox}{بحث کنید (روی تخته)}
بحث کنید که آیا در روش بازگشتی هم به پوزیشن انکدینگ نیاز داریم؟
\end{assignmentbox}

پاسخ کلاس: خیر. در \eng{RNN}/\eng{LSTM} عنصرها \emph{به‌ترتیب} پردازش می‌شوند و ترتیب در خودِ محاسبه نهفته است. مثلاً اگر حرف «ت» دو بار در یک دنبالهٔ نویسه‌ای بیاید، حالت پنهان هنگام دیدن «ت» دوم با «ت» اول فرق دارد، چون اطلاعات نویسه‌های قبلی را حمل می‌کند. در مقابل، در رویکرد موازی همهٔ عنصرها هم‌زمان پردازش می‌شوند و ترتیب خودبه‌خود وارد \eng{Attention} نمی‌شود.

\begin{table}[htbp]
\centering
\caption{منشأ اطلاعات ترتیب در دو رویکرد ساخت \eng{Embedding} بافتی.}
\begin{tabular}{lcc}
\toprule
 & ترتیبی (\eng{RNN}/\eng{LSTM}) & موازی (\eng{Self-Attention}) \\
\midrule
نحوهٔ پردازش & گام‌به‌گام & همهٔ عنصرها هم‌زمان \\
منشأ اطلاعات ترتیب & خودِ بازگشت (حالت پنهان) & باید صریحاً افزوده شود \\
نیاز به کدگذاری موقعیت & خیر (در بحث کلاس) & بله \\
\bottomrule
\end{tabular}
\end{table}

\section{سه پرسش طراحی}
مسئله در کلاس به سه پرسش شکسته شد:
\begin{enumerate}
  \item \textbf{محل}: کدگذاری موقعیت را در کدام بخش معماری قرار دهیم؟
  \item \textbf{نحوهٔ اعمال}: اطلاعات موقعیت را چگونه به \eng{Embedding} متصل کنیم؟ (جمع یا کانکت)
  \item \textbf{نحوهٔ مقداردهی}: مقدارها را چگونه تعیین کنیم؟ (یادگیری یا فرمول)
\end{enumerate}

\subsection{محل قرارگیری در معماری}
اطلاعات موقعیت باید \emph{بعد از \eng{Embedding} اولیه و پیش از ورود به محاسبات \eng{Attention}} اضافه شود، چون هنگام مقایسهٔ \eng{Query} و \eng{Key} مدل باید تفاوت جایگاه عنصرها را ببیند.

\begin{figure}[htbp]
\centering
\begin{tikzpicture}[rtl, 
  blk/.style={draw,rounded corners=3pt,minimum width=5.2cm,minimum height=0.85cm,align=center,font=\small},
  ar/.style={-{Stealth[length=2mm]},thick}]
  \node[blk,fill=gray!12,minimum width=7.6cm] (tok) at (0,0) {\RL{توکن‌های ورودی (کلمه، نویسه، واج، \lr{byte})}};
  \node[blk,fill=blue!10] (emb) at (0,1.6) {\lr{Embedding} اولیه: $E\in\mathbb{R}^{T\times d}$};
  \node[blk,fill=green!15,minimum width=2.8cm] (pe) at (4.6,3.2) {\RL{$P$: کدگذاری موقعیت}};
  \node[draw,circle,fill=white,inner sep=1pt] (plus) at (0,3.2) {$+$};
  \node[blk,fill=orange!15,minimum width=6.4cm] (att) at (0,4.8) {\lr{Attention}: $Q$، $K$، $V$ از $X=E+P$};
  \node[blk,fill=blue!10] (ctx) at (0,6.4) {\RL{\lr{Embedding} بافتی (پویا)}};
  \draw[ar] (tok)--(emb); \draw[ar] (emb)--(plus); \draw[ar] (plus)--(att); \draw[ar] (att)--(ctx);
  \draw[ar] (pe)--(plus);
\end{tikzpicture}
\caption{جای کدگذاری موقعیت در معماری: پس از \eng{Embedding} اولیه و پیش از \eng{Attention}. (شکل افزوده)}
\end{figure}

\begin{exambox}{انعطاف معماری}
کدگذاری موقعیت یک \emph{ماژول} است و محل و شیوهٔ استفاده از آن به معماری وابسته است. استاد اشاره کرد ممکن است در امتحان سؤال طراحی معماری بیاید، مثلاً اینکه اطلاعات موقعیت را در کدام بخش یک شبکه وارد می‌کنیم.
\end{exambox}

\subsection{نحوهٔ اعمال: کانکت یا جمع}
\textbf{کانکت:} بردار موقعیت جداگانه‌ای (مثلاً ۱۰ بعدی) به انتهای \eng{Embedding} ۵۱۲ بعدی چسبانده می‌شود ($512+10=522$)؛ مشکل: ابعاد نمایش بزرگ می‌شود، در حالی که ابعاد \eng{Embedding} یک تصمیم مهم طراحی است. \textbf{جمع:} بردار موقعیت هم‌اندازهٔ \eng{Embedding} است و ابعاد ثابت می‌ماند. این روش در کلاس انتخاب شد:
\[
\boxed{X \;=\; E + P}
\]

\subsection{ماتریس کدگذاری موقعیت}
اگر ابعاد \eng{Embedding} برابر ۵۱۲ باشد و موقعیت‌ها تا طول ۱۰۰ پوشش داده شوند:
\[
\boxed{P\in\mathbb{R}^{100\times 512}},\qquad X_{pos} = E_{token} + P_{pos}.
\]
هر \emph{سطر} یک موقعیت و هر \emph{ستون} یک بعد است؛ برای کلمهٔ موقعیت دهم، سطر دهم $P$ با \eng{Embedding} همان کلمه جمع می‌شود. (روی تخته ترانهادهٔ این نمایش کشیده شده بود؛ فقط قرارداد است.) مقدارهای $P$ را دو راه می‌توان ساخت: \emph{یادگرفتن} یا \emph{فرمول ثابت}.

\section{روش اول: کدگذاری موقعیت یادگرفتنی}
ماتریس $P$ مجموعه‌ای از پارامترهای قابل‌آموزش است (\eng{Learnable Position Embedding Matrix}) و مقدارهایش در آموزش یاد گرفته می‌شود، همان ایدهٔ \eng{Embedding} قابل‌آموزش در برابر \eng{One-Hot} با قاعدهٔ ثابت. در مثال کلاس $100\times512=51{,}200$ پارامتر اضافه می‌شود.

\paragraph{لایهٔ خطی نیست.} روی تخته: «اکتیویشن نداریم، ضرب ماتریسی نداریم (با لایهٔ خطی اشتباه نشود)». $P$ مستقیماً به \eng{Embedding} اضافه می‌شود؛ \eng{Bias}، فعال‌ساز یا ضرب ماتریسی ندارد، چون فقط می‌خواهیم برای هر موقعیت یک بردار داشته باشیم.

\begin{correctionbox}{دقت: جدول جست‌وجو}
در پیاده‌سازی، $P$ یک \term{جدول جست‌وجو}{Lookup Table} است (مثل ماتریس \eng{Embedding} کلمه‌ها): شمارهٔ موقعیت، سطر مربوط را برمی‌گرداند؛ از نظر ریاضی معادل ضرب \eng{One-Hot} موقعیت در $P$ است بدون اجرای واقعی ضرب.
\end{correctionbox}

\paragraph{آموزش با پس‌انتشار.} مثل بقیهٔ پارامترها: خروجی $X=E+P$ وارد بقیهٔ شبکه می‌شود، \eng{Loss} نسبت به خروجی نهایی و هدف همان مسئله حساب می‌شود (نه نسبت به خود \eng{Embedding})، و گرادیان با قاعدهٔ زنجیره‌ای به $E$ و $P$ می‌رسد. چون $X_t=E_{token_t}+P_t$، داریم $\dfrac{\partial L}{\partial P_t}=\dfrac{\partial L}{\partial X_t}$. نتیجه مهم: فقط سطرهایی از $P$ که در آموزش \emph{استفاده شده‌اند} گرادیان می‌گیرند.

\begin{figure}[htbp]
\centering
\begin{tikzpicture}[rtl, x=0.55cm,y=0.42cm]
  \foreach \r in {1,...,8}{\foreach \c in {1,...,8}{\draw[gray!50] (\c-1,-\r+1) rectangle (\c,-\r);}}
  \foreach \r in {1,2,3}{\foreach \c in {1,...,8}{\fill[green!35,opacity=.7] (\c-1,-\r+1) rectangle (\c,-\r);}}
  \foreach \c in {1,...,8}{\fill[gray!15] (\c-1,-7+1) rectangle (\c,-8);}
  \node[left] at (-0.2,-0.5) {\lr{pos 1}}; \node[left] at (-0.2,-1.5) {\lr{pos 2}}; \node[left] at (-0.2,-2.5) {\lr{pos 3}};
  \node[left] at (-0.2,-4.5) {$\vdots$}; \node[left] at (-0.2,-7.5) {\lr{pos 100}};
  \node[right,align=left,font=\small] at (8.4,-1.5) {سطرهای دیده‌شده: گرادیان غیرصفر\\و به‌روزرسانی مؤثر};
  \node[right,align=left,font=\small] at (8.4,-7.5) {سطرهای دیده‌نشده: گرادیان صفر\\(مقدار اولیه تقریباً بدون تغییر)};
\end{tikzpicture}
\caption{اگر دادهٔ آموزشی فقط توالی‌هایی تا طول ۳ داشته باشد، فقط سه سطر نخست $P$ گرادیان می‌گیرند. (شکل افزوده)}
\end{figure}

در حالت یادگرفتنی شرط صریحی برای متمایز بودن سطرهای $P$ نمی‌گذاریم و انتظار داریم داده، تابع خطا و ابعاد زیاد مدل را به نمایش مناسب برسانند (استاد گفت تضمین دقیق آن را می‌شود بیشتر بررسی کرد).

\section{روش دوم: کدگذاری موقعیت ثابت (فرمولی)}
نمایش موقعیت‌های مختلف باید \emph{قابل‌تمایز} باشد (مثلاً برای «شکارچی پلنگ را زد» در برابر «پلنگ شکارچی را زد»). هر خانهٔ ماتریس به دو چیز وابسته است: \emph{شمارهٔ موقعیت} ($pos$) و \emph{شمارهٔ بعد} ($j$)، یعنی $P(pos,j)=f(pos,j)$. روی تخته $PE(pos,j)=\sin\big(g(pos,j)\big)$ نوشته شد که $g$ زاویه‌ای ساخته‌شده از موقعیت و بعد است. چون $-1\le\sin(x)\le 1$ (و همین برای کسینوس)، مقدارها \emph{کراندار} هستند. استاد تأکید را روی منطق گذاشت، نه حفظ فرمول.

\begin{definitionbox}{فرمول استاندارد (\eng{Vaswani et al., 2017})}
برای $i=0,1,\dots,d/2-1$:
\[
PE(pos,\,2i)=\sin\!\Big(\frac{pos}{10000^{2i/d}}\Big),\qquad
PE(pos,\,2i+1)=\cos\!\Big(\frac{pos}{10000^{2i/d}}\Big).
\]
بعدها دو‌به‌دو (سینوس، کسینوس) با فرکانس $\omega_i=10000^{-2i/d}$ جفت می‌شوند؛ بعدهای کم‌شماره موج‌های سریع (تمایز موقعیت‌های نزدیک) و بعدهای پرشماره موج‌های کند (تمایز موقعیت‌های دور) می‌سازند.
\end{definitionbox}

\begin{correctionbox}{فرمول کلاس نسخهٔ ساده‌شده است}
فرمول $\sin(g(pos,j))$ کلاس نسخهٔ ساده‌شدهٔ فرمول استاندارد است (در نسخهٔ کامل بعدهای فرد کسینوس می‌گیرند). همچنین یادداشت «هیچ دو خانه‌ای زاویه یکسان ندارند» دقیق نیست (مثلاً با $d=512$ زاویهٔ ($pos=10$، بعد $i+64$) برابر زاویهٔ ($pos=1$، بعد $i$) است)؛ آنچه مهم است متفاوت‌بودن \emph{بردار کامل} هر موقعیت است.
\end{correctionbox}

\coursefigure{0.85}{ch10-pe-heatmap.png}{ماتریس $PE$ سینوسی استاندارد با $d=64$ و ۱۰۰ موقعیت (محور افقی: بعد، محور عمودی: موقعیت). بعدهای کم‌شماره (چپ) سریع نوسان می‌کنند و بعدهای پرشماره کندتر. (شکل افزوده)}

\begin{addedbox}{ویژگی‌های فرمول سینوسی}
\textbf{بدون پارامتر آموزشی}، و \textbf{قابل محاسبه برای هر موقعیت} (حتی دیده‌نشده در آموزش). همچنین $PE(pos+k)$ یک تبدیل خطی ثابت (دوران در هر جفت بعد) از $PE(pos)$ است و ضرب داخلی دو بردار موقعیت فقط به فاصلهٔ آن‌ها بستگی دارد.
\end{addedbox}

\section{مقایسهٔ دو روش}

\begin{table}[htbp]
\centering
\caption{مقایسهٔ کدگذاری موقعیت یادگرفتنی و ثابت.}
\begin{tabular}{p{3.2cm}p{5.2cm}p{5.2cm}}
\toprule
ویژگی & یادگرفتنی (\eng{Learned}) & ثابت (\eng{Fixed}) \\
\midrule
تعیین مقدارها & از طریق آموزش & با فرمول مشخص \\
پارامتر قابل‌آموزش & دارد ($100\times512$ در مثال) & ندارد \\
وابستگی به دادهٔ آموزشی & زیاد & کم \\
موقعیت‌های دیده‌نشده & ممکن است خوب یاد گرفته نشده باشند & با فرمول قابل محاسبه‌اند \\
چالش اصلی & پوشش موقعیت‌ها در دادهٔ آموزشی & انتخاب قاعدهٔ مناسب \\
\bottomrule
\end{tabular}
\end{table}

\section{سؤال مفهومی مهم استاد: دادهٔ آموزشی کوتاه}
\label{sec:ch10-exam}
دو مدل داریم: موقعیت در مدل اول با فرمول محاسبه می‌شود و در مدل دوم یاد گرفته می‌شود. اگر دادهٔ آموزشی عمدتاً جملهٔ کوتاه باشد، برای موقعیت‌های بلندتر کدام انتخاب بهتر است؟

\paragraph{پاسخ کلاس.} \emph{کدگذاری ثابت}. فرض کنید ماتریس یادگرفتنی ۱۰۰ سطر دارد ($P\in\mathbb{R}^{100\times 512}$) ولی دادهٔ آموزشی فقط توالی‌هایی تا طول ۳ دارد. گرادیان فقط به سطرهای ۱ تا ۳ می‌رسد و سطرهای مثلاً ۵۰ یا ۸۰ آموزش مؤثری نمی‌بینند؛ در آزمون با توالی به طول ۵۰ مدل از نمایش‌هایی استفاده می‌کند که هرگز یاد نگرفته و عملکرد افت می‌کند. در روش ثابت، بردار موقعیت ۵۰ با فرمول محاسبه می‌شود و لازم نیست در دادهٔ آموزشی دیده شده باشد.

\begin{exambox}{نکات روش‌شناختی این سؤال}
\begin{itemize}
  \item موضوع فقط \emph{تعداد} نمونه‌ها نیست، بلکه \emph{پوشش موقعیت‌ها} است: دادهٔ بسیار بزرگ ولی همگی کوتاه، تجربهٔ کافی دربارهٔ موقعیت‌های بالاتر نمی‌دهد؛ اگر دادهٔ آموزشی همهٔ موقعیت‌ها را پوشش دهد، روش یادگرفتنی هم خوب کار می‌کند.
  \item پاسخ «اندازهٔ ماتریس را کوچک بگیریم» فرض مسئله را عوض می‌کند: ظرفیت ۱۰۰ موقعیت است، دادهٔ آموزشی تا طول ۳، و اندازه تغییر نمی‌کند. استاد گفت در امتحان فرض‌ها را دقیق بیان خواهد کرد.
  \item \textbf{صورت پیشنهادی امتحان:} «دو مدل داریم که هر دو برای توالی‌هایی با حداکثر طول ۱۰۰ طراحی شده‌اند؛ مدل اول کدگذاری ثابت و مدل دوم یادگرفتنی دارد. اگر تمام دادهٔ آموزشی فقط شامل توالی‌هایی با طول حداکثر ۳ باشد، کدام روش برای موقعیت‌های بلندتر مناسب‌تر است و چرا؟»
\end{itemize}
\end{exambox}

\begin{correctionbox}{دقت: «ثابت بهتر است» تضمین کامل نیست}
مزیت روش ثابت این است که \emph{خودِ نمایش موقعیت} به آموزش هر سطر وابسته نیست؛ اما وزن‌های \eng{Attention} هم فقط روی موقعیت‌های کوتاه آموزش دیده‌اند و در عمل مدل‌های سینوسی هم روی طول‌های بلندتر از آموزش اغلب افت می‌کنند. پس روش ثابت \emph{مناسب‌تر} است، نه اینکه مشکل را کاملاً حل کند.
\end{correctionbox}

\section{کدگذاری موقعیت به توالی وابسته است، نه به «جمله»}
کدگذاری موقعیت ذاتاً به مفهوم «جمله» وابسته نیست. عنصرهای توالی می‌توانند کلمه، نویسه، واج یا بایت باشند و کدگذاری موقعیت فقط مشخص می‌کند هر عنصر \emph{چندمین عنصر توالی ورودی} است: در ورودی نویسه‌ای موقعیت پنجم یعنی نویسهٔ پنجم، و در ورودی کلمه‌ای یعنی کلمهٔ پنجم. این به طراحی معماری و \eng{Tokenization} بستگی دارد.

\section{طول زمینه (\eng{Context Length})}
اگر توالی ۱۱۰ عنصر باشد ولی ماتریس موقعیت فقط ۱۰۰ موقعیت داشته باشد چه می‌شود؟
\begin{itemize}
  \item حداکثر طول ورودی یک تصمیم طراحی است؛ برای توالی‌های بلندتر می‌توان ورودی را به بخش‌های کوچک‌تر شکست و از پنجره (\eng{Context Window}) استفاده کرد.
  \item اندازهٔ پوشش موقعیت به مدل و \eng{Task} بستگی دارد؛ مثلاً مدل نویسه‌ای معمولاً به موقعیت‌های بیشتری از مدل کلمه‌ای نیاز دارد.
  \item در روش یادگرفتنی موقعیت ۱۰۱ سطری ندارد؛ در روش فرمولی می‌توان $PE(101)$ را حساب کرد، ولی کیفیت آن تضمین‌شده نیست.
\end{itemize}

\section{مثال عددی کامل: \eng{Attention} در رمزگشا}
\label{sec:ch10-example}
استاد روی تخته مثال عددی دست‌نویسی از \eng{Attention} در معماری رمزگذار--رمزگشا (همان \eng{G2P}) حل کرد: ورودی رمزگذار و رمزگشا هر دو \eng{Embedding} نویسه (\eng{CE}) به‌علاوهٔ \eng{PE} هستند و \eng{Query} رمزگشا با \eng{Key}/\eng{Value} رمزگذار بردار بافت $C$ را می‌سازد.

\paragraph{گام ۱: \eng{Key} و \eng{Value}.} $E\in\mathbb{R}^{3\times 2}$ (سه توکن) با ضرب در $W_K$ و $W_V$ (هر دو $2\times2$) به $K$ و $V$ ($3\times2$) تبدیل می‌شود:

\begin{figure}[htbp]
\centering
\resizebox{0.95\textwidth}{!}{%
\begin{tikzpicture}[rtl, ar/.style={-{Stealth[length=2.5mm]},thick,mainblue},
  fr/.style={draw,minimum width=1.2cm,minimum height=1.1cm}]
  \node (E) at (0,0) {$\underset{3\times2}{E=\begin{bmatrix}1&2\\3&4\\5&6\end{bmatrix}}$};
  \node (K) at (9,1.9) {$\underset{3\times2}{K=\begin{bmatrix}9&10\\11&12\\13&14\end{bmatrix}}$};
  \node (V) at (9,-1.9) {$\underset{3\times2}{V=\begin{bmatrix}15&16\\17&18\\19&20\end{bmatrix}}$};
  \draw[ar] (1.9,0.5)--node[above,sloped]{$\times\,\underset{2\times2}{W_K}$} (6.6,1.9);
  \draw[ar] (1.9,-0.5)--node[above,sloped]{$\times\,\underset{2\times2}{W_V}$} (6.6,-1.9);
  \node[draw=red!70!black,circle,thick,inner sep=2pt] at (-2.4,0) {۱};
\end{tikzpicture}}
\caption{گام ۱: $E$ ($3\times2$) در $W_K$ و $W_V$ ($2\times2$) ضرب می‌شود و $K$ و $V$ ($3\times2$) به دست می‌آیند. روی تخته $W_K$ و $W_V$ فقط به‌صورت قاب‌های خالی کشیده شده بودند و مقدارشان نوشته نشده بود. (بازسازی‌شده از تختهٔ کلاس)}
\end{figure}

\paragraph{گام ۲: \eng{Query}.} توکن رمزگشا پس از جمع با کدگذاری موقعیت در $W_Q$ ضرب می‌شود: $Q=(E_{dec}+P_{pos})\,W_Q\in\mathbb{R}^{1\times 2}$؛ یعنی موقعیت پیش از ساخت $Q$ وارد نمایش می‌شود.

\paragraph{گام ۳: امتیازها.} $\text{\lr{scores}} = Q K^{\top}\in\mathbb{R}^{1\times 3}$ ، ابعاد $(1\times2)\cdot(2\times3)$.

\begin{figure}[htbp]
\centering
\resizebox{0.95\textwidth}{!}{%
\begin{tikzpicture}[rtl]
  \node[draw=red!70!black,circle,thick,inner sep=2pt] at (-0.5,0) {۲};
  \node[draw=orange,rounded corners=3pt,fill=orange!12,inner sep=5pt] (ep) at (2.5,0) {$\underset{1\times2}{(E_{dec}+PE)}$};
  \node at (5.1,0) {$\times$};
  \node[draw,minimum width=1.3cm,minimum height=1.1cm,label=below:{\small \lr{$W_Q$ ($2\times2$)}}] at (6.5,0) {};
  \node at (8.0,0) {$=$};
  \node (Q) at (10.3,0) {$\underset{1\times2}{Q=\begin{bmatrix}1&2\end{bmatrix}}$};
  \node[draw=red!70!black,circle,thick,inner sep=2pt] at (-0.5,-3) {۳};
  \node (sc) at (7.2,-3) {$\underset{1\times2}{Q}\;\cdot\;\underset{2\times3}{K^{\top}=\begin{bmatrix}9&11&13\\10&12&14\end{bmatrix}}\;=\;\underset{1\times3}{\begin{bmatrix}29&35&41\end{bmatrix}}$};
  \node[font=\small,text=gray] at (7.2,-4.3) {\RL{امتیازها} \lr{scores}};
\end{tikzpicture}}
\caption{گام‌های ۲ و ۳: \eng{Embedding} توکن رمزگشا با \eng{PE} جمع و در $W_Q$ ضرب می‌شود تا $Q$ ساخته شود (بالا)؛ سپس $QK^{\top}$ امتیازها را می‌دهد (پایین). مقدارهای $Q$ و امتیازها دست‌خط به‌سختی خوانایی بودند و در اینجا بر پایهٔ خوانش توضیح‌داده‌شده در متن آمده‌اند. (بازسازی‌شده از تختهٔ کلاس)}
\end{figure}

\begin{correctionbox}{نکتهٔ خوانایی و \eng{Softmax}}
اعداد دست‌نویس $Q$ و امتیازها روی تخته ناخوانا بود؛ خوانش سازگار $Q=(1,2)$ و $QK^\top=(29,35,41)$ است (منطق محاسبه از مقدار دقیق آن‌ها مستقل است). همچنین امتیاز خام ضریب توجه نیست: پیش از ضرب در $V$ باید با \eng{Softmax} (معمولاً پس از تقسیم بر $\sqrt{d_k}$) به وزن‌هایی با مجموع ۱ تبدیل شود. ضرب مستقیم امتیازهای خام در $V$ مقدار بی‌معنی $(1809,1914)$ می‌دهد.
\end{correctionbox}

\paragraph{گام ۴: بردار بافت.} $C=\sum_{i=1}^{3}\alpha_i V_i\in\mathbb{R}^{1\times 2}$. با امتیازهای بالا
\[
\mathrm{softmax}(29,35,41)\approx(0{,}000006;\,0{,}0025;\,0{,}9975),\qquad C\approx(18{,}995;\;19{,}995)
\]
(توجه «تیز»؛ تقسیم بر $\sqrt{d_k}$ آن را نرم‌تر می‌کند).

\paragraph{گام ۵: لجیت.} $\text{\lr{logits}} = C\,W + b$ با $W\in\mathbb{R}^{2\times 5}$ و $b\in\mathbb{R}^{1\times 5}$ (واژگان مثال ۵ عضو دارد)، و سپس \eng{Softmax} روی واژگان احتمال هر توکن بعدی را می‌دهد.

\paragraph{ماسک علّی.} در رمزگشا $QK^{\top}$ برای همهٔ موقعیت‌ها محاسبه و مثلث بالا‌راست (توکن‌های آینده) با $-\infty$ پیش از \eng{Softmax} ماسک می‌شود تا وزن آن‌ها صفر شود؛ هر توکن فقط به خودش و توکن‌های قبلی نگاه می‌کند:
\[
\text{\lr{mask}}=\begin{bmatrix}0&-\infty&-\infty\\0&0&-\infty\\0&0&0\end{bmatrix},\qquad
\alpha=\mathrm{softmax}\big(QK^\top/\sqrt{d_k}+\text{\lr{mask}}\big).
\]

\begin{exambox}{یادداشت‌های پایانی روی تخته}
روی تخته دو یادداشت آمده بود: همان پرسش بحث دربارهٔ نیاز روش بازگشتی به کدگذاری موقعیت (بخش \ref{sec:ch10-rnn}) و «مدل‌های اتورگرسیو (\eng{Autoregressive}) و \eng{Teacher Forcing} در زمان آموزش» که هنوز تدریس نشد و موضوع جلسهٔ بعد است.
\end{exambox}

\section{جمع‌بندی}
\begin{itemize}
  \item \eng{Self-Attention} موازی بدون اطلاعات موقعیت نسبت به ترتیب بی‌اطلاع است؛ \eng{RNN} ترتیب را ضمناً از بازگشت می‌گیرد.
  \item کدگذاری موقعیت پس از \eng{Embedding} اولیه و پیش از \eng{Attention} با $X=E+P$ اعمال می‌شود؛ جمع ابعاد را ثابت نگه می‌دارد.
  \item $P$ را می‌توان یاد گرفت (جدول جست‌وجو، فقط سطرهای دیده‌شده گرادیان می‌گیرند) یا با فرمول سینوسی ساخت (بدون پارامتر، برای هر موقعیت قابل محاسبه).
  \item موقعیت به توالی ورودی مربوط است نه لزوماً «جمله»؛ حداکثر طول یک تصمیم طراحی است.
\end{itemize}

\begin{exambox}{نکات مهم برای مرور و امتحان}
\begin{enumerate}
  \item چرا \eng{Self-Attention} بدون کدگذاری موقعیت دو جملهٔ «ایران آمریکا را شکست داد» و «آمریکا ایران را شکست داد» را اشتباه می‌گیرد؟ چرا \eng{RNN} این مشکل را ندارد؟
  \item محل کدگذاری موقعیت (پس از \eng{Embedding} و پیش از \eng{Attention}) و دلیل آن.
  \item کانکت در برابر جمع؛ $X=E+P$ و $P\in\mathbb{R}^{100\times512}$؛ تفسیر سطر و ستون.
  \item یادگرفتنی: پارامتر قابل‌آموزش، \eng{Backprop}، بدون فعال‌ساز و ضرب ماتریسی (لایهٔ خطی نیست).
  \item ثابت: ورودی‌های فرمول ($pos$ و $j$)، کران $[-1,1]$ سینوس و کسینوس.
  \item سؤال امتحانی مهم: دادهٔ آموزشی کوتاه (طول ۳) و ماتریس ۱۰۰ موقعیتی؛ چرا کدگذاری ثابت مناسب‌تر است؟ فرض‌ها را دقیق بنویسید.
  \item ممکن است سؤال طراحی بیاید: اطلاعات موقعیت را در کجای یک شبکهٔ داده‌شده وارد می‌کنید؟
  \item مثال عددی \eng{Attention}: ابعاد هر گام ($E,K,V:3\times2$، $Q:1\times2$، $QK^\top:1\times3$، $C:1\times2$، لجیت $1\times5$).
\end{enumerate}
\end{exambox}
